\documentclass[10pt,a4paper]{amsart}
\usepackage{amsmath,amssymb,amsthm,graphicx}
% For commutative diagrams we use Xy-pic
\usepackage[all]{xy}

% We use 2cell for 2-commutative diagrams.
\xyoption{2cell}
% Theorem environments.
%
\theoremstyle{plain}
\newtheorem{theorem}[subsection]{Theorem}
\newtheorem{proposition}[subsection]{Proposition}
\newtheorem{lemma}[subsection]{Lemma}

\theoremstyle{definition}
\newtheorem{definition}[subsection]{Definition}
\newtheorem{example}[subsection]{Example}
\newtheorem{exercise}[subsection]{Exercise}
\newtheorem{situation}[subsection]{Situation}

\theoremstyle{remark}
\newtheorem{remark}[subsection]{Remark}
\newtheorem{remarks}[subsection]{Remarks}

\numberwithin{equation}{subsection}

% Macros
%
\def\lim{\mathop{\mathrm{lim}}\nolimits}
\def\colim{\mathop{\mathrm{colim}}\nolimits}
\def\Spec{\mathop{\mathrm{Spec}}}
\def\Hom{\mathop{\mathrm{Hom}}\nolimits}
\def\Ext{\mathop{\mathrm{Ext}}\nolimits}
\def\SheafHom{\mathop{\mathcal{H}\!\mathit{om}}\nolimits}
\def\SheafExt{\mathop{\mathcal{E}\!\mathit{xt}}\nolimits}
\def\Sch{\mathit{Sch}}
\def\Mor{\mathop{\mathrm{Mor}}\nolimits}
\def\Ob{\mathop{\mathrm{Ob}}\nolimits}
\def\Sh{\mathop{\mathit{Sh}}\nolimits}
\def\NL{\mathop{N\!L}\nolimits}
\def\CH{\mathop{\mathrm{CH}}\nolimits}
\def\proetale{{pro\text{-}\acute{e}tale}}
\def\etale{{\acute{e}tale}}
\def\QCoh{\mathit{QCoh}}
\def\Ker{\mathop{\mathrm{Ker}}}
\def\Im{\mathop{\mathrm{Im}}}
\def\Coker{\mathop{\mathrm{Coker}}}
\def\Coim{\mathop{\mathrm{Coim}}}

% Boxtimes
%
\DeclareMathSymbol{\boxtimes}{\mathbin}{AMSa}{"02}

%
% Macros for moduli stacks/spaces
%
\def\QCohstack{\mathcal{QC}\!\mathit{oh}}
\def\Cohstack{\mathcal{C}\!\mathit{oh}}
\def\Spacesstack{\mathcal{S}\!\mathit{paces}}
\def\Quotfunctor{\mathrm{Quot}}
\def\Hilbfunctor{\mathrm{Hilb}}
\def\Curvesstack{\mathcal{C}\!\mathit{urves}}
\def\Polarizedstack{\mathcal{P}\!\mathit{olarized}}
\def\Complexesstack{\mathcal{C}\!\mathit{omplexes}}
% \Pic is the operator that assigns to X its picard group, usage \Pic(X)
% \Picardstack_{X/B} denotes the Picard stack of X over B
% \Picardfunctor_{X/B} denotes the Picard functor of X over B
\def\Pic{\mathop{\mathrm{Pic}}\nolimits}
\def\Picardstack{\mathcal{P}\!\mathit{ic}}
\def\Picardfunctor{\mathrm{Pic}}
\def\Deformationcategory{\mathcal{D}\!\mathit{ef}}
\title{Every element coprime to b up to a unit makes a monogenic extension of a ring containing a unit of the form a - by}
\author{The Stacks Project Authors}
\date{stacks-project:etale-cohomology.tex:local-snapshot-20261011; GNU FDL 1.2 or later}
\begin{document}
\maketitle
\noindent\textit{Editorial note.} All delivered mathematical text is the author's own, copied from the pinned Stacks Project chapter file etale-cohomology.tex; no statement or proof step was rewritten, reordered or removed.\par\medskip
\noindent\textit{Editorial note.} The counted proof is the author's complete argument: The counted proof is the author's complete argument, including the stepwise construction of the coefficients a\_1, ..., a\_n and the closing prime-ideal contradiction.\par\medskip
\noindent\textit{Editorial note.} The label lemma-find-extension defines the statement; the counted proof is the proof environment immediately following the lemma in the source.\par\medskip
\noindent\textit{Editorial note.} The preamble is the chapter's own native preamble excerpt already used for this source: the Xy-pic setup (native preamble lines 14-18) and the native theorem declarations with the mathematical macros (native preamble lines 154-217). The class block and the externaldocument lines of the native preamble are not included, because the assembler supplies the wrapper class and those external documents are not part of this material.\par\medskip
\noindent\textit{Allowed standard prerequisites (editorial).}\begin{itemize}
\item Monogenic extension of a ring. Same chapter, the definition of a monogenic extension A subset B used throughout this section. Supplies the target object B = A[y]/(f) of the construction and the meaning of the conclusion.
\item Root of unity modulo an ideal. Same chapter (and standard usage), a in A such that a\textasciicircum{}n = 1 mod bA for some n >= 1. Supplies the exponent n whose powers of a are units modulo b\textasciicircum{}m A, and hence modulo b.
\end{itemize}\medskip
\begin{lemma}
\label{lemma-find-extension}
Let $A$ be a ring. Let $a, b \in A$ such that
$aA + bA = A$ and $a \bmod bA$ is a root of unity.
Then there exists a monogenic extension $A \subset B$
and an element $y \in B$ such that $u = a - by$ is a unit.
\end{lemma}

\begin{proof}
Say $a^n \equiv 1 \bmod bA$. In particular $a^i$ is a unit modulo
$b^mA$ for all $i, m \geq 1$. We claim there exist
$a_1, \ldots, a_n \in A$ such that
$$
1 = a^n + a_1 a^{n - 1}b + a_2 a^{n - 2}b^2 + \ldots + a_n b^n
$$
Namely, since $1 - a^n \in bA$ we can find an element $a_1 \in A$
such that $1 - a^n - a_1 a^{n - 1} b \in b^2A$ using the unit property
of $a^{n - 1}$ modulo $bA$. Next, we can find an element $a_2 \in A$
such that $1 - a^n - a_1 a^{n - 1} b - a_2 a^{n - 2} b^2 \in b^3A$.
And so on. Eventually we find $a_1, \ldots, a_{n - 1} \in A$
such that $1 - (a^n + a_1 a^{n - 1}b + a_2 a^{n - 2}b^2 +
\ldots + a_{n - 1} ab^{n - 1}) \in b^nA$. This allows us to find
$a_n \in A$ such that the displayed equality holds.

\medskip\noindent
With $a_1, \ldots, a_n$ as above we claim that setting
$$
B = A[y]/(y^n + a_1 y^{n - 1} + a_2 y^{n - 2} + \ldots + a_n)
$$
works. Namely, suppose that $\mathfrak q \subset B$ is a prime ideal
lying over $\mathfrak p \subset A$. To get a contradiction
assume $u = a - by$ is in $\mathfrak q$. If $b \in \mathfrak p$
then $a \not \in \mathfrak p$ as $aA + bA = A$ and hence $u$ is not
in $\mathfrak q$. Thus we may assume $b \not \in \mathfrak p$, i.e.,
$b \not \in \mathfrak q$. This implies that $y \bmod \mathfrak q$
is equal to $a/b \bmod \mathfrak q$. However, then we obtain
$$
0 = y^n + a_1 y^{n - 1} + a_2 y^{n - 2} + \ldots + a_n =
b^{-n}(a^n + a_1 a^{n - 1}b + a_2 a^{n - 2}b^2 + \ldots + a_nb^n) =
b^{-n}
$$
a contradiction. This finishes the proof.
\end{proof}

\end{document}
